% ════════════════════════════════════════════════════════════════════
% SECTION 2 — BACKGROUND AND RELATED WORK
% ════════════════════════════════════════════════════════════════════
\section{Background and Related Work}\label{sec:related}

\subsection{Solar flare prediction from photospheric magnetograms}

The physical premise of data-driven flare forecasting is that the
photospheric magnetic field of an active region encodes the stress
budget that a future flare will release. Schrijver~\cite{schrijver2007}
demonstrated that a single scalar---the R-value, the unsigned flux
within high-gradient polarity-separation corridors---separates
major-flaring regions from non-flaring ones with striking reliability,
and this proxy remains among the strongest single flare predictors
identified to date. Bobra and Couvidat~\cite{bobra2015} operationalised
the machine-learning view of the problem, showing that a support-vector
machine trained on 25~SHARP parameters derived from SDO/HMI vector
magnetograms can forecast M-class flares with useful skill, and that
feature selection matters as much as model choice. Subsequent
community-wide blind tests, most notably the multi-method comparison of
Leka et al.~\cite{leka2019}, established two durable findings: (i)~no
single method dominates across all event definitions and skill
metrics, and (ii)~the effective ceiling on skill is set by the quality
and homogeneity of the training data rather than by classifier
complexity.

The European FLARECAST project~\cite{georgoulis2021} pushed the
operational agenda further, combining multiple forecasting methods and
data sources in an open framework, and it explicitly confronted the
data-heterogeneity problem that motivates the present paper. On the
data side, the SWAN-SF benchmark of Angryk et
al.~\cite{angryk2020} released a multivariate time-series dataset of
4{,}098 active-region (MVTS) tracks with 24 physical parameters per
timestep and flare-history labels, specifically designed to make flare
prediction reproducible and comparable; it has since become the
community's reference benchmark, and a family of cleaned and rebalanced
variants (random undersampling, Tomek-link cleaning, and TimeGAN-based
synthetic augmentation \cite{yoon2019}) has been curated to make the
benchmark tractable for supervised learning. Recent deep-learning work
on SWAN-SF continues to raise the centralised ceiling: contrastive
pretraining and multifaceted preprocessing \cite{eskandari2024} and
LSTM-style temporal architectures \cite{hassani2025} exploit the
sequential structure of the 60-minute observation windows. All of these
works, without exception, assume that the full training corpus is
available in one place.

\subsection{Federated learning under statistical heterogeneity}

Federated learning was introduced by McMahan et
al.~\cite{mcmahan2017} in the FedAvg formulation: clients perform
several epochs of local SGD on private data, and a server aggregates
the resulting weights by sample-size-weighted averaging. The defining
difficulty, identified almost immediately, is statistical
heterogeneity: when client data distributions differ---the non-IID
setting---local optima diverge from the global optimum, and naive
averaging can oscillate or diverge \cite{zhao2018}. FedProx
\cite{li2020} addresses this by adding a proximal term
$(\mu/2)\norm{\w-\w^{(t)}}^{2}$ to each client's local objective that
anchors local training to the current global model, tolerating partial
work and heterogeneous local computation. SCAFFOLD \cite{karimireddy2020}
attacks the same problem with control variates that correct client
drift using variance reduction. The comprehensive survey of Kairouz et
al.~\cite{kairouz2021} situates these algorithms within the broader
open-problem landscape, including robust aggregation, personalisation,
and privacy. Cross-silo FL---few, reliable, institutionally owned
clients---is the deployment regime relevant to space-weather agencies,
and it contrasts with the cross-device regime in its stability
requirements and long-lived participants.

A distinct and, for this study, decisive strand concerns normalisation
statistics under heterogeneity. FedBN~\cite{li2021fedbn} established
that BatchNorm running statistics are a per-client feature that
weight-averaged federation must \emph{not} naively transport, and that
keeping them local outperforms both FedAvg and FedProx under feature
shift. Wang et al.~\cite{wang2023bn} supplied the convergence analysis
of exactly this mismatch, Guerraoui et al.~\cite{guerraoui2024bn}
showed corrected statistics restore centralised-level evaluation
performance, and BN-SCAFFOLD~\cite{bnscaffold2024} documents BN drift
crippling SCAFFOLD specifically. The parameters-only transport bug we
diagnose in Section~\ref{sec:res-bn} is an engineering-level
instantiation of this literature's central mechanism --- the mechanism
is textbook; the composite failure chain (initialised-buffer evaluation
times a saturated monitor times a checkpoint rule) is the
pipeline-specific part.

\subsection{Class imbalance in federated learning}

Flare prediction inherits an extreme class-imbalance regime: in
SWAN-SF's natural test distribution, roughly 1.9\% of windows precede
M/X-class flares, while the rebalanced training partitions sit near
parity. The focal loss of Lin et al.~\cite{lin2017}, originally
developed for dense object detection, down-weights easy examples with a
$(1-p_t)^{\gamma}$ modulating factor and concentrates gradient on hard
positives; it has become the default imbalance workhorse. Sarkar et
al.~\cite{sarkar2020} adapted focal loss specifically to FL in the
Fed-Focal formulation, in which the class-weighting parameter is
adjusted using information available at the server about the global
class distribution, so that clients with different local positivity
rates do not implicitly reweight the global objective. In parallel,
synthetic oversampling remains a standard pre-processing response:
SMOTE \cite{chawla2002} interpolates minority-class neighbours, and
mixup \cite{zhang2018} creates convex combinations of training pairs.
Under federation, rebalancing must be applied locally per client, since
the pooled distribution is precisely what no participant can inspect.

\subsection{Positioning of this work and search protocol}
\label{sec:litsearch}

Two research threads---physics-grounded flare prediction and
heterogeneity-robust federated learning---have developed largely
independently, and one prior work joins them: Fu, Wan, Huang, Han, and
Peng posted ``Federated Transfer Learning for Soalr Flare Forecasting''
[title sic] to Research Square on 10~July~2023, formatted for
\emph{Solar Physics}~\cite{fu2023}. That preprint applies federated
transfer learning to flare forecasting on its own data substrate; it
remains unrefereed, and we found no evidence of formal journal
publication as of this revision. It does not evaluate SWAN-SF, reports
no provenance audit, and, to our knowledge, no follow-up has appeared.
We note that its evaluation protocol---a random 7:1.5:1.5 split of
pooled windows (AUC 0.936 on the Tianchi flare dataset)---is precisely
the random-split convention whose instance-sharing risk this paper
quantifies for SWAN-SF; the comparison is offered as context, not as a
score comparison across datasets.
Because novelty claims are factual statements about the literature, we
document the search protocol behind ours rather than relying on ``to
the best of our knowledge''. The search was last
repeated on 29~September 2026 across general web, arXiv, NASA ADS,
and publisher indexes (ScienceDirect, IEEE) reached through a web
search engine, using four query strings: (i)~\emph{federated learning
solar flare prediction}; (ii)~\emph{federated learning space weather
forecasting}; (iii)~\emph{federated learning SWAN-SF benchmark
magnetogram}; (iv)~\emph{cross-silo federated learning geophysics
heliophysics distributed observatories}. Records were included if
they apply federated or distributed multi-client learning without raw
data pooling to flare prediction or space-weather forecasting; records
were excluded if they concern centralised machine learning for flares
(the case for all entries of query~i), federated forecasting of
\emph{solar power or irradiance}, federated \emph{terrestrial}
weather forecasting, or federated computing on spacecraft without a
scientific forecasting task (the categories returned by query~ii).
Screening the 30 records returned by the four queries yielded
\emph{zero} in-scope records; the raw result listings are committed
with the repository (\texttt{docs/lit\_search/}, each listing
embedding its query date since v4.6). That screening
\emph{missed} the Fu et al.\ preprint above---exactly the
grey-literature coverage limitation this section flags---and the
preprint surfaced only in a targeted verification pass during the
present revision. We note two further limitations of this protocol:
web-engine coverage of preprint servers is incomplete, and the
screening was single-reviewer, so a PRISMA-style dual-screened log
remains queued repository work. Three facts bound the novelty claim honestly. The leakage
phenomenon is documented prior art: the SWAN-SF creators warned that
temporal coherence invalidates randomised
sampling~\cite{angryk2019}, Ahmadzadeh et al.\ repeated the
warning~\cite{ahmadzadeh2021}, and the cleaned export's release paper
prescribes temporally-preceding folds~\cite{eskandari2024}; the
BatchNorm-transport mechanism is textbook since
FedBN~\cite{li2021fedbn}. Within those limits, the present work is
the first federated evaluation on the SWAN-SF benchmark (in both
its cleaned and raw substrates, against centralised baselines of
matched architecture), the first \emph{instance-level} provenance
audit of the cleaned export (normalisation-invariant matching of
every window to the raw benchmark), and the first leakage-free
re-evaluation of federated flare prediction under the benchmark's
intended protocol. It also contributes a rarely documented artifact:
an honest engineering log of the failure modes encountered when FL
meets a severely imbalanced physical-science benchmark
(Appendix~\ref{sec:appendix}).
